\section{L'entreprise et le périmètre audité (G1)}

\subsection{Organisation de MEDICA-SERVICES}

MEDICA-SERVICES emploie 220 personnes sur quatre implantations et dispose d'une DSI de 14 personnes \piece{(P01, P02)}. La figure~\ref{fig:organigramme} reprend son organisation. Seuls deux liens sont établis par les pièces. La DG approuve la stratégie SSI que la DSI applique, et le RSSI est rattaché à la DSI. Les autres rattachements ne sont pas précisés par \piece{P02} et ne sont donc pas tracés.

\begin{figure}[H]
\centering
\resizebox{\linewidth}{!}{% Organigramme de MEDICA-SERVICES (P01, P02, P09). Seuls les liens établis par les pièces sont tracés.
\begin{tikzpicture}[font=\footnotesize,
  dir/.style={draw=accent, fill=accentclair, rounded corners=2pt, align=center, inner sep=5pt, text width=#1},
  dir/.default=3.2cm,
  fonc/.style={draw=gris, fill=white, rounded corners=2pt, align=center, inner sep=4pt, text width=#1},
  fonc/.default=2.9cm,
  ext/.style={draw=gris, dashed, fill=grisclair, rounded corners=2pt, align=center, inner sep=4pt, text width=#1},
  ext/.default=2.9cm,
  lab/.style={font=\scriptsize, fill=white, inner sep=1.5pt, align=center}]

  % Direction
  \node[dir=4.6cm, fill=accent, text=white] (dg) at (0,0)
    {\textbf{Direction générale (DG)}\\[1pt]\scriptsize autorité du périmètre, approuve la stratégie SSI et le budget};

  % Fonctions citées par P02
  \node[dir] (dsi) at (-5.2,-2.6) {\textbf{DSI}\\\scriptsize 14 personnes\\ applique la stratégie SSI};
  \node[fonc] (do) at (-1.6,-2.6) {\textbf{Direction des opérations}\\\scriptsize DO};
  \node[fonc] (crh) at (1.75,-2.6) {\textbf{Contrats et RH}\\\scriptsize contrats clients et prestataires};
  \node[ext] (dpo) at (5.1,-2.6) {\textbf{DPO}\\\scriptsize externalisé};

  \node[dir=3.2cm, fill=white] (rssi) at (-5.2,-5.0) {\textbf{RSSI}\\\scriptsize 1 personne à mi-temps};

  % Responsables opérationnels (rattachement non précisé)
  \node[fonc=2.4cm] (sup) at (-1.3,-5.0) {\textbf{Resp. support}};
  \node[fonc=2.4cm] (inf) at (1.45,-5.0) {\textbf{Resp. infrastructures}};
  \node[fonc=2.4cm] (dev) at (4.2,-5.0) {\textbf{Resp. développement}};
  \node[fonc=8.0cm] (ast) at (1.45,-6.35) {\textbf{24 techniciens d'astreinte}\\\scriptsize sauvegardes, déploiements et reprise (P01, P14, P15)};
  \begin{scope}[on background layer]
    \node[draw=filet, dashed, rounded corners=3pt, fit=(sup)(inf)(dev)(ast), inner sep=6pt,
      label={[font=\scriptsize\itshape, text=gris]above:responsables opérationnels, rattachement non précisé par P02}] {};
  \end{scope}

  % Liens établis
  \draw[accent, thick, dashed, -{Stealth[length=2mm]}] (dg.south) -- ++(0,-0.55) -| node[lab, pos=0.27, below=1pt]{stratégie SSI approuvée par la DG, appliquée par la DSI (P02)} (dsi.north);
  \draw[accent, very thick] (dsi.south) -- node[lab, right=2pt]{rattachement\\ hiérarchique (P02)} (rssi.north);
  \draw[gris, dotted, very thick, -{Stealth[length=2mm]}] (rssi.west) -- ++(-0.45,0) |- node[lab, pos=0.72, above=1pt]{compte rendu oral, 2 fois par an (P09)} (dg.west);

  % Parties externes et sites
  \node[ext=4.1cm] (cli) at (-4.4,-7.9) {\textbf{12 établissements de soins}\\\scriptsize clients, reprise sous 8 h (P01)};
  \node[ext=4.1cm] (sec) at (0.15,-7.9) {\textbf{Prestataire du site de secours}\\\scriptsize réplication sous contrat (P03, T01)};
  \node[ext=4.1cm] (sit) at (4.7,-7.9) {\textbf{Sites et effectif}\\\scriptsize Toulouse, Bordeaux, Montpellier, Montpellier-Sud, 220 personnes (P01)};

  % Légende
  \node[anchor=north west, font=\scriptsize, align=left] at (3.6,0.7) {%
    \tikz\draw[accent, very thick] (0,0)--(0.55,0); rattachement établi\\[1pt]
    \tikz\draw[accent, thick, dashed] (0,0)--(0.55,0); lien fonctionnel\\[1pt]
    \tikz\draw[gris, dotted, very thick] (0,0)--(0.55,0); reporting\\[1pt]
    \tikz\draw[gris, dashed, fill=grisclair] (0,0) rectangle (0.55,0.18); externe};
\end{tikzpicture}
}
\caption{Organigramme de MEDICA-SERVICES d'après P01, P02 et P09.}
\label{fig:organigramme}
\end{figure}

\begin{table}[H]
\centering\small
\caption{Qui gère quoi en matière de SSI}
\begin{tabularx}{\linewidth}{@{}L{3.1cm} Y L{1.6cm} L{3.4cm}@{}}
\toprule
\entete \hd{Acteur} & \hd{Rôle établi par les pièces} & \hd{Pièces} & \hd{Rôle dans notre plan} \\
\midrule
Direction générale & Autorité du périmètre, approuve la stratégie SSI et l'enveloppe de 120\,000\,€ & P02, P09, P16 & Décide, accepte les risques résiduels \\
DSI & Applique la stratégie SSI. Infrastructure et comptes par hypothèse & P02 & Pilote A08 et les règles FW \\
RSSI (mi-temps) & Recense quelques risques avec la DSI, rend compte oralement deux fois par an & P02, P09, P10 & Pilote A01, A02, A03 et l'investigation \\
Resp. infrastructures & Cité par P02. Sauvegardes et correctifs par hypothèse & P02 & Pilote A05 et A09 \\
Support et astreinte & Connaît la procédure de restauration, 24 techniciens d'astreinte & P01, P15 & Exécute les tests de restauration \\
Contrats et RH & Gestion des contrats et des ressources humaines & P02 & Clauses prestataires, sensibilisation \\
DPO externalisé & Fonction externalisée & P02 & Consulté sur les données de santé \\
\bottomrule
\end{tabularx}
\end{table}

\subsection{Périmètre retenu}

La méthode recommande un périmètre homogène en termes de responsabilité. Nous retenons le service d'hébergement et d'exploitation rendu aux clients, placé sous l'autorité de la DG.

\begin{table}[H]
\centering\small
\caption{Périmètre retenu}
\begin{tabularx}{\linewidth}{@{}L{3.1cm} Y L{2.1cm}@{}}
\toprule
\entete \hd{Élément} & \hd{Contenu retenu} & \hd{Pièces} \\
\midrule
Bénéficiaires & 12 établissements de soins, reprise sous 8 h des applications critiques & P01 \\
Actifs essentiels & MEDI-PORTAL et sa base, identifiants patients et données de santé, ERP, ticketing & P03, P05, T01 \\
Services supports & 35 serveurs ou VM, Active Directory, Microsoft 365, VPN et bastion, sauvegardes, salle principale et salle de secours & P03, T01, T09 \\
Interconnexions & Portail HTTPS public, liaisons contractuelles avec les clients, réplication vers le secours, SaaS & P03, P07, T01 \\
Exclusions & SI des clients, fonctionnement interne des SaaS et du prestataire de secours, réseaux de Bordeaux, Montpellier et Montpellier-Sud non décrits, tout test actif & P01, T01, T04 \\
\bottomrule
\end{tabularx}
\end{table}

\subsection{Limites et informations manquantes}

Les pièces sont des fragments. Une absence de justificatif est traitée comme une incertitude à lever et non comme la preuve qu'une action n'a pas eu lieu. Trois manques pèsent le plus sur nos conclusions.

\begin{table}[H]
\centering\small
\caption{Informations manquantes et demandes de preuves}
\begin{tabularx}{\linewidth}{@{}L{4.0cm} Y Y@{}}
\toprule
\entete \hd{Information manquante} & \hd{Conséquence} & \hd{Demande de preuve} \\
\midrule
RTO et RPO par service, dépendances \piece{(T01, T09)} & Engagement de 8 h invérifiable & Analyse d'impact, annexes contractuelles, entretien avec le responsable infrastructures \\
Liste exhaustive des comptes privilégiés et accès conditionnel \piece{(T04)} & Taux MFA valable pour l'export seul & Export complet annuaire et SaaS, registre des exceptions \\
Suite donnée à AL-1048 et intégrité des journaux \piece{(T07, T08)} & Événement du 9 octobre non qualifiable & Journaux bruts scellés, historique du ticketing, entretien avec le titulaire d'admin-old \\
\bottomrule
\end{tabularx}
\end{table}
